\documentclass[
  aps,
  pra,
  reprint,            % two-column, journal look
  superscriptaddress, % for author affiliations
  amsmath,amssymb,
  notitlepage
]{revtex4-2}

% Encoding & fonts
\usepackage[utf8]{inputenc}
\usepackage[T1]{fontenc}

% General useful packages
\usepackage[hypertexnames=false]{hyperref} % hyperlinks
\usepackage{url}            % simple URL typesetting
\usepackage{booktabs}       % professional-quality tables
\usepackage{amsfonts}       % blackboard math symbols
\usepackage{nicefrac}       % compact symbols for 1/2, etc.
\usepackage{microtype}      % microtypography

\usepackage{graphicx}       % graphics
\graphicspath{{media/}}     % organize your images under media/ folder

\usepackage[normalem]{ulem}
\usepackage{tikz}
\usetikzlibrary{positioning,quotes}
\usepackage{xcolor}
\usepackage{soul}
\usepackage{mathtools}
\usepackage{quantikz}
\usepackage{cancel}

\usepackage{algorithm}
\usepackage{algpseudocode}

%\usepackage{comment}
\newcommand{\fm}[1]{\textcolor{blue}{[Felix: #1]}}
\newcommand{\mm}[1]{\textcolor{blue}{[Matthias: #1]}}
\newcommand{\bb}[1]{\textcolor{red}{[Bora: #1]}}


\begin{document}


\author{Bora Baran}
\email{b.baran@fz-juelich.de}
\affiliation{Forschungszentrum Jülich, Institute of Quantum Control,%
 Peter Grünberg Institut (PGI-8), 52425 Jülich, Germany}
\affiliation{Institute for Theoretical Physics, University of Cologne, Zuelpicher Strasse 77, 50937 Cologne, Germany}

\author{Tommaso Calarco}
\affiliation{Forschungszentrum Jülich, Institute of Quantum Control,%
 Peter Grünberg Institut (PGI-8), 52425 Jülich, Germany}
\affiliation{Institute for Theoretical Physics, University of Cologne, Zuelpicher Strasse 77, 50937 Cologne, Germany}
\affiliation{Dipartimento di Fisica e Astronomia, Università di Bologna, 40127 Bologna, Italy}

\author{Matthias M. Müller}
%\email{ma.mueller@fz-juelich.de}
\affiliation{Forschungszentrum Jülich, Institute of Quantum Control,%
 Peter Grünberg Institut (PGI-8), 52425 Jülich, Germany}

\author{Felix Motzoi}
%\email{f.motzoi@fz-juelich.de}
\affiliation{Forschungszentrum Jülich, Institute of Quantum Control,%
 Peter Grünberg Institut (PGI-8), 52425 Jülich, Germany}
\affiliation{Institute for Theoretical Physics, University of Cologne, Zuelpicher Strasse 77, 50937 Cologne, Germany}


\title{Noise-Robust Single-Pulse Synthesis of Tripartite Entangling Interactions\\
in a Realistic NV-Center Spin Register}

\begin{abstract}
\noindent
Fast multi-qubit entangling operations are commonly compiled from sequences of one- and two-qubit gates, increasing execution time and accumulated control error. Here we develop noise-robust single-pulse synthesis of tripartite entangling interactions in a realistic nitrogen-vacancy (NV) center register formed by an electronic spin and two nearby $^{13}\mathrm C$ nuclei. Our support-selective objective resolves the diagonal-frame phase generator into one-body, pairwise, and tripartite contributions. It fixes the requested three-body coefficient, suppresses unwanted pairwise terms, and leaves correctable local coefficients open. This construction is a direct extension of local-equivalence and perfect-entangler optimization from two-qubit nonlocal classes to selected many-body interaction supports. We use the remaining local freedom to reduce a trajectory-dependent electronic dephasing-exposure functional while retaining the target tripartite content. The previously obtained $11$-component Fourier pulses for $\exp[i(\pi/4)Z_AZ_BZ_C]$ and $\exp[i(\pi/4)X_AZ_BZ_C]$ serve as high-fidelity warm starts. Robustness is evaluated under electronic Ornstein--Uhlenbeck detuning on the Cartesian grid $T_{2,A}^{\ast}\in\{2,5,10,20,50\}\,\mu\mathrm{s}$ and $\tau_c\in\{0.1,0.3,1,3,15,30\}\,\mu\mathrm{s}$, with three independently characterized diamond environments overlaid as source-consistent points. Time-resolved electron-manifold population and $Z_A$-variance integrals connect the optimized trajectories to the stochastic gate results. \bb{Insert final robust-pulse fidelities, exposure reductions, and experimental-scenario losses after the production runs.}
\end{abstract}


\maketitle

%\tableofcontents

\section*{Introduction}

Quantum optimal control provides systematic methods for shaping external fields to steer a quantum system through a desired transformation, enabling precise implementation of logical gates, state transfers, and error-resilient operations~\cite{koch2022quantum}. It has become a central tool for engineering high-performance operations across a broad range of quantum hardware platforms, including trapped ions~\cite{monz201114}, neutral-atom arrays~\cite{evered2023high,levine2019parallel,omran2019generation}, superconducting circuits~\cite{li2024hardware,zhao2025microwave,barnes2017fast,li2024experimental,gao2025ultrafast}, and solid-state spin systems such as nitrogen-vacancy centers and related platforms~\cite{dong2024time,ChouNVOptim,wang2025circuit}. These developments are supported by a mature algorithmic framework, including gradient-based and variational optimization methods such as GRAPE~\cite{khaneja2005optimal,de2011second,motzoi2011optimal,hogben2011spinach}, Krotov methods~\cite{konnov1999deterministic,Reich07062014,goerz2019krotov,reich2012monotonically,goerz2014optimal}, and dCRAB or chopped-basis parametrizations~\cite{caneva2011chopped,rach2015dressing,doria2011optimal,muellercontrol}. In many gate-synthesis settings, these methods optimize a cost functional defined by a target unitary, so that the realized operation approximates one prescribed matrix representative, typically up to a global phase.

\medskip
\noindent
This full-unitary viewpoint is natural when the complete action of a gate is operationally relevant. In many multi-qubit tasks, however, only selected nonlocal content is required, while local operations can be corrected, commuted into neighboring circuit layers, or absorbed into frame updates. Local-invariant and local-equivalence-class methods exploit this freedom by optimizing two-qubit gates through their nonlocal content rather than one fixed matrix representative. Closely related work optimizes directly over perfect-entangler classes, thereby enlarging the admissible control landscape without changing the desired entangling capability~\cite{zhang2003geometric,rezakhani2004characterization,muller2011optimizing,goerz2015optimizing,watts2015optimizing,frey2026optimal}.

\medskip
\noindent
The present construction is a simple extension of this principle from two-qubit local-equivalence classes to selected many-body interaction supports. In a local frame where the requested interaction is diagonal, a Walsh--Hadamard transform separates the computational-basis phases into one-body, pairwise, and tripartite Pauli-string coefficients. We constrain the requested tripartite coefficient and suppress unwanted pairwise coefficients, while leaving the one-body sector open. The resulting family is not identical to the full quotient by arbitrary local unitaries; rather, it is the support-resolved analogue appropriate to the diagonal-frame interaction considered here.

\medskip
\noindent
The unused local freedom also provides a control resource. Among propagators with the same selected three-body content, the optimizer can choose trajectories that expose the electronic actuator less strongly to longitudinal detuning. We therefore augment the terminal support-selective objective with a differentiable time-integrated electronic exposure term. The high-fidelity pulses obtained without this term are retained as reference solutions and used as warm starts, allowing a controlled comparison between reference and exposure-penalized gates at fixed duration, Fourier basis size, and target interaction content.

NV centers combine optically addressable electronic spin states with hyperfine-coupled nuclear spins, enabling initialization, readout, and coherent control of multi-spin registers at or near room temperature~\cite{phila,Schir,neumann,rembold2020introduction}. Nearby $^{13}\mathrm C$ nuclear spins can serve as long-lived register qubits, while the NV electronic spin provides an optically readable actuator and mediator for conditional control~\cite{Taminiau,childress,wang2025circuit}. This register architecture is therefore a natural setting in which to test whether useful higher-order interaction content can be generated directly by a shaped microwave pulse rather than compiled from a longer sequence of lower-order gates.

We consider a logical three-qubit subsystem formed by the NV electronic spin, denoted by $A$, and two proximal $^{13}\mathrm C$ nuclear spins, denoted by $B$ and $C$. We synthesize two target operations,
\begin{equation}
U_{ZZZ}=e^{i\frac{\pi}{4}Z_AZ_BZ_C},
\qquad
U_{XZZ}=e^{i\frac{\pi}{4}X_AZ_BZ_C}.
\label{eq:introduction_target_gates}
\end{equation}
For both operations, the terminal cost targets the tripartite phase coefficient and suppresses all pairwise coefficients. The $ZZZ$ operation is treated directly in the computational basis, whereas the $XZZ$ operation is optimized in a local Hadamard frame on the electronic qubit. The resulting gates are implemented by single microwave pulses and subsequently tested under explicit ensembles of temporally correlated dephasing trajectories.

\section{System of Interest: NV-Center Spin Register}
\label{sec:system_of_interest}

\subsection{NV-center spin register}
\label{subsec:nv_register}

The physical system considered in this work is a negatively charged nitrogen-vacancy center in diamond coupled to nearby nuclear spins. An NV center is formed when one carbon atom in the diamond lattice is replaced by a nitrogen impurity and an adjacent lattice site is vacant. Among its possible charge states, the negatively charged NV$^-$ center is of particular interest because of its spin-dependent optical transitions and magneto-optical properties~\cite{phila,Schir}. In the following, we refer to the NV$^-$ charge state simply as the NV center.

The electronic ground state of the NV center is a spin triplet. Two of its ground-state spin levels, conventionally chosen from the $\ket{0_e}$ and $\ket{-1_e}$ manifold under an applied magnetic field, form an effective electronic qubit. Optical excitation allows spin initialization and spin-dependent fluorescence readout, while microwave fields enable coherent manipulation of the electronic transition~\cite{phila}. These features have made NV centers a versatile platform for quantum sensing, quantum information processing, and solid-state spin control.

A central resource for multi-qubit operation is the hyperfine coupling between the NV electronic spin and surrounding nuclear spins. Every NV center is intrinsically coupled to the nitrogen nuclear spin that defines the defect. In addition, nearby $^{13}\mathrm{C}$ nuclei may be present in the diamond lattice. Since the lattice is mostly occupied by spinless $^{12}\mathrm{C}$, while $^{13}\mathrm{C}$ has natural abundance of about $1.1\%$ and nuclear spin $1/2$, only a subset of lattice sites host nuclear-spin qubits~\cite{neumann}. When a $^{13}\mathrm{C}$ nucleus is sufficiently close to the NV center, its hyperfine interaction produces a resolvable splitting of the electronic transition. The splitting depends sensitively on the nuclear position in the lattice, allowing different nearby nuclear spins to be spectrally distinguished and, in favorable cases, individually addressed~\cite{neumann,nizovtsev2018nonflipping,nizovtsev2021hyperfine}.

The resulting NV-nuclear system naturally forms a small spin register. Nuclear spins can exhibit coherence times much longer than the electronic spin and can therefore act as memory qubits, while the electronic spin can be used for initialization, readout, and conditional control of the register~\cite{Taminiau,childress}. Register spins can be read out by mapping their state onto the NV electronic spin, for example through entangling operations between the electronic and nuclear degrees of freedom~\cite{Rong2015,Wei2015}. Initialization and readout fidelities can be affected by charge-state fluctuations and optical pumping dynamics~\cite{taminiaudetection,neumannsingle}, but the platform remains highly attractive because coherent multi-spin control can be achieved under experimentally accessible conditions.

In the simulations below, the logical register consists of the NV electronic spin, denoted by $A$, coupled to two proximal $^{13}\mathrm{C}$ nuclear spins, denoted by $B$ and $C$. The intrinsic $^{14}\mathrm{N}$ nuclear spin is prepared in a fixed nuclear-spin manifold and acts as a spectator degree of freedom. The logical subsystem $A\otimes B\otimes C$ is therefore the three-qubit register on which we synthesize tripartite entangling operations.

\subsection{Effective driven Hamiltonian and register parameters}
\label{sec:spin_model}

The dynamics of the driven NV-center spin register are described by an effective Hamiltonian obtained from the laboratory-frame NV Hamiltonian after a two-level reduction of the electronic spin, a secular approximation for the hyperfine interaction, and a rotating-wave approximation at the driven electronic transition. A full derivation is given in Appendix~\ref{app:nv_model}. In the interaction picture with respect to the free precessing terms, the Hamiltonian used for the control simulations is
\begin{equation}
\begin{aligned}
&\notag H_{\mathrm{NV}}(t)\\
&=\frac{\Omega(t)}{2}\sum_{\mathbf m}
\big[\sigma_x\cos((\omega_{\mathrm{MW}}-\Lambda(\mathbf m))\,t)\\
&\hspace{1.0cm}-\sigma_y\sin((\omega_{\mathrm{MW}}-\Lambda(\mathbf m))\,t)\big] \otimes \ket{\mathbf m}\!\bra{\mathbf m} \\
&-\;\frac{\Omega(t)}{2}\big[\sigma_x\sin((\omega_{\mathrm{MW}}-\Lambda_s) t)+\sigma_y\cos((\omega_{\mathrm{MW}}-\Lambda_s) t)\big] \\
&\hspace{1.0cm}\otimes\sum_i \theta_i \big[I'_{ix}\cos(\delta_i t-\phi_i)+I'_{iy}\sin(\delta_i t-\phi_i)\big].
\end{aligned}
\label{eq:H_NV_RWA_main}
\end{equation}
The first term describes coherent microwave driving of the electronic qubit conditioned on the nuclear-spin configuration. The second term captures weak microwave-mediated coupling to nearby nuclear spins that arises from the relative tilt between the nuclear quantization axes in the electronic $\ket{0_e}$ and $\ket{-1_e}$ manifolds. Here, $\Omega(t)$ is the time-dependent Rabi frequency controlled by the applied microwave field with carrier frequency $\omega_{\mathrm{MW}}$. The quantity
\begin{equation}
\Lambda_s = D - \gamma_e B_0
\end{equation}
is the bare electronic transition frequency between the selected ground-state levels, including the zero-field splitting $D$ and the Zeeman shift due to the static magnetic field $B_0$. The configuration-dependent transition frequency $\Lambda(\mathbf m)$ depends on the nuclear-spin configuration $\mathbf m=(m_N,\{m_{C_i}\})$, where $m_N$ and $m_{C_i}$ label the nuclear spin projections of the $^{14}\mathrm{N}$ and $^{13}\mathrm{C}$ nuclei. The quantities $\delta_i=(\gamma_i B_0+\omega_i)/2$ describe weak effective transverse nuclear modulations, while $\theta_i$ and $\phi_i$ characterize the local hyperfine-axis misalignment of each nuclear spin. The static field and electronic parameters used throughout are
\begin{equation}
\begin{aligned}
B_0&=0.45\,\mathrm{T},
&\frac{D}{2\pi}&=2.87\,\mathrm{GHz},\\
\frac{\gamma_e}{2\pi}&=28.024\,\mathrm{GHz/T}.
\end{aligned}
\label{eq:electronic_register_parameters}
\end{equation}
The nuclear-spin parameters defining the simulated register are
collected in Table~\ref{tab:nv_register_parameters}. The first carbon
represents the strongly coupled register class used by Chen
\emph{et al.}~\cite{Chen}. The second carbon is associated with the
near-stable St1 coupling family identified in the lattice-resolved
calculations of Nizovtsev \emph{et al.}; the calculated longitudinal
coupling $A_{ZZ}/2\pi\simeq-1.011\,\mathrm{MHz}$ and weak
nonsecular scale $A_{\mathrm{nd}}/2\pi\simeq14.5\,\mathrm{kHz}$ agree
with the parameters used here to the quoted precision
~\cite{nizovtsev2018nonflipping}. Related calculations show how
selected $^{13}\mathrm C$ lattice-site families can provide
NV--nuclear register spins with predictable hyperfine properties
~\cite{nizovtsev2021hyperfine}.

\begin{table}[t]
\centering
\caption{Parameters defining the simulated
NV--$^{14}\mathrm N$--$^{13}\mathrm C_1$--$^{13}\mathrm C_2$
register. Frequencies are quoted after division by $2\pi$.
$^{13}\mathrm C_1$ represents the strongly coupled register class
used in Ref.~\cite{Chen}; $^{13}\mathrm C_2$ matches the weakly
transverse, near-stable St1 lattice family calculated in
Refs.~\cite{nizovtsev2018nonflipping,nizovtsev2021hyperfine}.}
\label{tab:nv_register_parameters}
\scriptsize
\setlength{\tabcolsep}{3.0pt}
\renewcommand{\arraystretch}{1.10}
\begin{tabular}{@{}lcccc@{}}
\toprule
Register spin and class
& $\gamma/2\pi$
& $A_{zz}/2\pi$
& $A_{\perp}/2\pi$
& $Q/2\pi$ \\
& (MHz/T) & (MHz) & (MHz) & (MHz) \\
\midrule
$^{14}\mathrm N$ (intrinsic), $I=1$
& $3.077$ & $-2.14$ & $0$ & $-5.01$ \\
$^{13}\mathrm C_1$ (strong), $I=\tfrac12$
& $10.71$ & $2.281$ & $0.240$ & $0$ \\
$^{13}\mathrm C_2$ (St1-like), $I=\tfrac12$
& $10.71$ & $-1.011$ & $0.014$ & $0$ \\
\bottomrule
\end{tabular}
\end{table}

The two carbon spins therefore define a realistic representative
register assembled from experimentally and microscopically supported
coupling classes. We do not claim that both parameter sets were
measured simultaneously in one particular NV center or that the
St1-like coupling identifies one unique lattice coordinate.

For a given control pulse, the driven evolution defines the propagator
\begin{equation}
U_{\mathrm{pulse}}(\mathbf p,T)
=
\mathcal{T}
\exp\!\left(
-i\int_0^T
H_{\mathrm{NV}}(\Omega(t;\mathbf p),t)\,dt
\right),
\label{eq:U_pulse_definition}
\end{equation}
where $T$ is the pulse duration and $\mathbf p$ denotes the pulse parameters. The logical three-qubit propagator used in the target evaluation is the restriction of Eq.~\eqref{eq:U_pulse_definition} to the $A\otimes B\otimes C$ subspace with the $^{14}\mathrm{N}$ spin fixed in the chosen spectator manifold; we denote this projected propagator by $U^{A\otimes B\otimes C}_{\mathrm{pulse}}(\mathbf p,T)$.

\subsection{Pulse parametrization}
\label{sec:pulse_parametrization}

We control the system by optimizing the temporal shape of the Rabi frequency appearing in Eq.~\eqref{eq:H_NV_RWA_main}. The pulse is represented by a compact analytic parametrization,
\begin{equation}
\Omega(t;\mathbf p) = \chi(t)\sum_{i=1}^{N_c} a_{i} \cos(\omega_{c,i} t + \phi_{c,i}),
\label{eq:pulse_parametrization}
\end{equation}
where $a_i$ are frequency amplitudes, $\omega_{c,i}$ are control frequencies, and $\phi_{c,i}$ are control phases. The optimization variables are collected into
\begin{equation}
\mathbf p =
(a_{1},\ldots,a_{N_c},\;
\omega_{c,1},\ldots,\omega_{c,N_c},\;
\phi_{c,1},\ldots,\phi_{c,N_c}).
\label{eq:pulse_parametrization_example}
\end{equation}
The multiplicative envelope $\chi(t)$ is a smooth turn-on and turn-off window. We use a tapered-cosine window with taper fraction $\alpha$,
\begin{equation}
\chi(t)=\tfrac12\left[1-\cos\!\left(2\pi\frac{t}{T}\right)\right], 
\qquad t\in[0,\alpha T],
\end{equation}
and set $\alpha=0.15$ in the numerical examples. The optimized pulse parameters are provided in the public data repository~\cite{baran2026_multiqubit_control}. Both tripartite gates reported below use $N_c=11$, so that they are compared within the same direct Fourier control family.

This type of frequency-domain ansatz is closely related to chopped-basis and CRAB-style parametrizations~\cite{caneva2011chopped,rach2015dressing,doria2011optimal,muellercontrol}, while retaining compatibility with gradient-based or hybrid optimization strategies used in quantum-control applications~\cite{khaneja2005optimal,de2011second,motzoi2011optimal,sorensen2018quantum,sorensen2020optimization,preti2022continuous}. It provides a compact and experimentally interpretable representation of the microwave control field.

\section{Support-Selective Control Objective}
\label{sec:control_objective}

The control objective is designed to generate the requested tripartite interaction without unnecessarily fixing all local phases of the logical operation. Let
\begin{equation}
U_L(\mathbf p,T)
=
U_{\mathrm{pulse}}^{A\otimes B\otimes C}(\mathbf p,T)
\label{eq:logical_pulse_short}
\end{equation}
denote the logical three-qubit propagator obtained from the full NV-register evolution. We first choose a local frame $R$ in which the desired Pauli-string interaction is diagonal,
\begin{equation}
U_F(\mathbf p,T)
=
R^\dagger U_L(\mathbf p,T)R.
\label{eq:framed_logical_propagator}
\end{equation}
For the $ZZZ$ target, $R=I$. For the $XZZ$ target, we take
\begin{equation}
R=H_A\otimes I_B\otimes I_C,
\end{equation}
for which
\begin{equation}
R^\dagger X_AZ_BZ_C R=Z_AZ_BZ_C.
\end{equation}
Thus, the same diagonal-frame cost can be used for both physical interactions.

Suppose first that $U_F$ is diagonal in the computational basis. We write
\begin{equation}
U_F=e^{-i\Phi},
\qquad
\Phi=
\sum_{x\in\{0,1\}^3}
\varphi_x\ket{x}\!\bra{x},
\label{eq:diagonal_phase_generator}
\end{equation}
where a consistent branch is chosen for the phases $\varphi_x$. Every diagonal Hermitian operator on three qubits has the expansion
\begin{equation}
\Phi
=
\sum_{S\subseteq\{A,B,C\}}
\phi(S)Z_S,
\qquad
Z_S=\prod_{j\in S}Z_j,
\label{eq:diagonal_phase_expansion}
\end{equation}
with $Z_\emptyset=I$. Because the diagonal Pauli strings are orthogonal,
\begin{equation}
\operatorname{Tr}(Z_SZ_{S'})=8\,\delta_{S,S'},
\end{equation}
the coefficient on support $S$ is
\begin{equation}
\boxed{
\phi(S)
=
\frac{1}{8}\operatorname{Tr}(\Phi Z_S)
=
\frac{1}{8}
\sum_{x\in\{0,1\}^3}
(-1)^{\sum_{j\in S}x_j}\varphi_x
}
\label{eq:support_phase_coordinate}
\end{equation}
where $x_j$ is the bit associated with qubit $j$. Equation~\eqref{eq:support_phase_coordinate} is the Walsh--Hadamard transform of the computational-basis phase map. It directly separates the global, one-body, pairwise, and tripartite coefficients of the diagonal phase generator. Explicit three-qubit expressions are listed in Appendix~\ref{app:explicit_3q_invariants}.

For the two gates considered here, the desired diagonal-frame coefficients are
\begin{equation}
\phi^\star(ABC)=\frac{\pi}{4},
\qquad
\phi^\star(AB)=
\phi^\star(AC)=
\phi^\star(BC)=0.
\label{eq:target_phase_coordinates}
\end{equation}
The one-body coefficients $\phi(A)$, $\phi(B)$, and $\phi(C)$ are not included in the terminal target. They represent local rotations in the selected frame and are removed after optimization by a fixed local correction.

We compare the selected coefficients with their target values through the periodic cost
\begin{equation}
\begin{aligned}
\mathcal J_{\mathrm{phase}}
&=
w_3
\left[
1-\cos\!\left(2[\phi(ABC)-\pi/4]\right)
\right]
\\
&\quad
+w_2
\sum_{S\in\{AB,AC,BC\}}
\left[
1-\cos\!\left(2\phi(S)\right)
\right].
\end{aligned}
\label{eq:support_selective_phase_cost}
\end{equation}
The factor of two implements the natural $\pi$ periodicity of every nonidentity Pauli-string coefficient, since
\begin{equation}
e^{-i(\theta+\pi)Z_S}=-e^{-i\theta Z_S}
\end{equation}
for nonempty $S$, and the minus sign is only a global phase.

During optimization, the framed propagator need not be exactly diagonal. We therefore extract $\varphi_x$ from the diagonal overlaps $\bra{x}U_F\ket{x}$ and add the diagonality penalty
\begin{equation}
\mathcal J_{\mathrm{diag}}
=
1-\frac{1}{8}
\sum_{x\in\{0,1\}^3}
\left|\bra{x}U_F\ket{x}\right|^2.
\label{eq:diagonality_penalty}
\end{equation}
For a unitary logical propagator, $\mathcal J_{\mathrm{diag}}=0$ exactly when every computational-basis state is preserved up to a phase in the selected frame.

To use the open one-body sector for dephasing robustness, we propagate the eight logical computational-basis states under the noiseless control. Let
\begin{equation}
\ket{\psi_x(t)}=U(t)V\ket{x},
\qquad x\in\{0,1\}^3,
\end{equation}
where $V$ embeds the logical subspace in the complete register Hilbert space. We define the normalized electronic dephasing-exposure functional
\begin{equation}
\boxed{
\mathcal J_{\mathrm e}
=
\frac{1}{8T}
\sum_{x\in\{0,1\}^3}
\int_0^T
\left[
1-
\left\langle\psi_x(t)\middle|\widetilde Z_A\middle|\psi_x(t)\right\rangle^2
\right]dt .
}
\label{eq:electron_dephasing_exposure}
\end{equation}
The integrand is the variance of $\widetilde Z_A$ for the propagated basis state and vanishes whenever the trajectory remains in one electronic $Z$ manifold. It is therefore a direct state-trajectory measure of susceptibility to longitudinal detuning. By contrast, the raw population of one electron manifold averaged over a complete logical basis is fixed by unitarity and cannot serve as an optimization objective.

For visualization, we additionally track the population transferred to the electronic manifold opposite to that of each input state. If $z_x=\langle x|V^\dagger\widetilde Z_AV|x\rangle=\pm1$, this population is
\begin{equation}
p_{\mathrm{opp}}^{(x)}(t)
=
\frac{1-z_x\langle\psi_x(t)|\widetilde Z_A|\psi_x(t)\rangle}{2},
\label{eq:opposite_manifold_population}
\end{equation}
and its logical-basis and time average is denoted $\overline{I}_{p}$. This population integral is a trajectory diagnostic; the stochastic OU propagation remains the decisive robustness test.

The complete numerical objective is
\begin{equation}
\mathcal J
=
\lambda_F\mathcal J_{F}
+\lambda_{\mathrm{phase}}\mathcal J_{\mathrm{phase}}
+\lambda_{\mathrm{diag}}\mathcal J_{\mathrm{diag}}
+\lambda_{\mathrm{leak}}\mathcal J_{\mathrm{leak}}
+\lambda_{\mathrm e}\mathcal J_{\mathrm e}
+\lambda_{\mathrm{reg}}\mathcal J_{\mathrm{reg}},
\label{eq:complete_nv_cost}
\end{equation}
where $\mathcal J_F$ is the local-corrected terminal infidelity, $\mathcal J_{\mathrm{leak}}$ penalizes population outside the logical subspace, and $\mathcal J_{\mathrm{reg}}$ contains weak pulse-shape regularization. The support weights remain $w_3=0.5$ and $w_2=0.2$. Starting from the reference solutions, we increase $\lambda_{\mathrm e}$ by continuation and select the lowest-exposure pulse that satisfies the prescribed terminal-fidelity threshold. All local corrections are calibrated once from the resulting noiseless pulse and remain fixed in the noise ensembles.

After optimization, the first-order coefficients determine the correction
\begin{equation}
U_{\mathrm{loc}}^{(F)}
=
\exp\!\left[
-i\sum_{j\in\{A,B,C\}}\phi(j)Z_j
\right]
\label{eq:fixed_local_correction_frame}
\end{equation}
in the diagonal frame. For the $ZZZ$ gate this is already the computational-frame correction. For the $XZZ$ gate it is transformed back with $R$, producing an $X_A$ correction and $Z_B,Z_C$ corrections. The correction is calibrated once from the noiseless optimized pulse and is kept fixed throughout the stochastic simulations.

\section{Experimentally Anchored Electron-Spin Dephasing Model}
\label{sec:ou_dephasing_model}

To assess coherence loss during the shaped gates, we propagate the complete driven register under explicit classical detuning trajectories. The primary analysis is restricted to longitudinal frequency fluctuations of the NV electronic qubit. This isolates dephasing of the directly driven actuator spin and avoids introducing additional nuclear-noise parameters that are not independently characterized for the selected register. Nuclear-spin noise can be included by the same construction, but it is not part of the results reported here.

Let
\begin{equation}
\widetilde Z_A
=
\left(\ket{0_e}\!\bra{0_e}-\ket{-1_e}\!\bra{-1_e}\right)
\otimes I_{^{14}\mathrm N}
\otimes I_{^{13}\mathrm C_1}
\otimes I_{^{13}\mathrm C_2}
\label{eq:embedded_electron_z}
\end{equation}
denote the electronic Pauli-$Z$ operator embedded in the full simulated register. For realization $r$, the Hamiltonian propagated during the gate is
\begin{equation}
\boxed{
H_r(t)
=
H_{\mathrm{NV}}(t)
+
\frac{\beta_r(t)}{2}\widetilde Z_A .
}
\label{eq:noisy_full_hamiltonian}
\end{equation}
As throughout the Hamiltonian model, $\beta_r(t)$ and its stationary width are expressed in angular-frequency units. Positive and negative values of $\beta_r(t)$ shift the electronic transition above and below the nominal microwave resonance. For the nominal free Hamiltonian used to obtain Eq.~\eqref{eq:H_NV_RWA_main}, $\widetilde Z_A$ commutes with the electron-diagonal free terms and therefore retains the same operator form in the interaction picture. No echo, dynamical-decoupling, or additional refocusing pulse is inserted into the gate propagation; the only applied control is the optimized waveform $\Omega(t)$.

The detuning is represented by a stationary Ornstein--Uhlenbeck process~\cite{UhlenbeckOrnstein1930},
\begin{equation}
\mathrm d\beta_r(t)
=
-\frac{\beta_r(t)}{\tau_c}\,\mathrm dt
+
\sqrt{\frac{2\sigma_A^2}{\tau_c}}\,\mathrm dW_r(t),
\label{eq:ou_sde}
\end{equation}
with zero mean, variance $\sigma_A^2$, and autocorrelation
\begin{equation}
\left\langle\beta_r(t)\beta_r(t')\right\rangle
=
\sigma_A^2e^{-|t-t'|/\tau_c}.
\label{eq:ou_correlation}
\end{equation}
The correlation time $\tau_c$ controls how rapidly the detuning changes. For $\tau_c$ comparable to a gate duration, it evolves appreciably during the shaped pulse; for $\tau_c\gg T$, it is almost constant within one realization while varying between realizations.

For a simulation step $\Delta t$, the process is initialized in its stationary distribution, $\beta_{r,0}\sim\mathcal N(0,\sigma_A^2)$, and propagated with the exact finite-step transition~\cite{gillespie1996ou},
\begin{equation}
\beta_{r,n+1}
=
e^{-\Delta t/\tau_c}\beta_{r,n}
+
\sigma_A\sqrt{1-e^{-2\Delta t/\tau_c}}\,\xi_{r,n},
\label{eq:ou_exact_discrete}
\end{equation}
where the $\xi_{r,n}$ are independent standard normal variables. This update preserves the prescribed stationary variance and exponential autocorrelation on the chosen propagation grid.

The Hamiltonian is treated as piecewise constant over each interval $[t_n,t_{n+1})$, with the control and stochastic detuning evaluated on the same grid. The short-time propagator is therefore
\begin{equation}
U_{r,n}
=
\exp\!\left\{
-i\left[
H_{\mathrm{NV}}(t_n)
+
\frac{\beta_{r,n}}{2}\widetilde Z_A
\right]\Delta t
\right\},
\label{eq:noisy_timestep_propagator}
\end{equation}
and the complete realization-dependent gate is
\begin{equation}
\boxed{
U_r(T)=U_{r,N-1}\cdots U_{r,1}U_{r,0}.
}
\label{eq:noisy_ordered_product}
\end{equation}
Because the driven Hamiltonian generally does not commute with $\widetilde Z_A$, the noise changes the driven trajectory itself---including rotation axes, transition amplitudes, and conditional phases---rather than acting as a random $Z$ rotation appended after an ideal gate. Numerical convergence is checked by repeating the material study at a finer propagation grid.

For stationary Gaussian longitudinal noise, averaging the accumulated random phase gives~\cite{cywinski2008dephasing}
\begin{equation}
W_R(t)=\exp[-\chi_R(t)],
\qquad
\chi_R(t)
=
\sigma_A^2\tau_c^2
\left(
\frac{t}{\tau_c}-1+e^{-t/\tau_c}
\right).
\label{eq:ou_ramsey_envelope}
\end{equation}
When a Ramsey time is used as the source-native characterization, the condition $W_R(T_{2,A}^{\ast})=e^{-1}$ determines
\begin{equation}
\sigma_A
=
\frac{1}{
\tau_c
\sqrt{
T_{2,A}^{\ast}/\tau_c
-1
+e^{-T_{2,A}^{\ast}/\tau_c}
}
}.
\label{eq:ou_sigma_ramsey}
\end{equation}
The characterization experiments cited below employ Ramsey, Hahn-echo, or multipulse spectroscopy to infer environmental properties. These characterization pulses are not transferred to our gate simulation: only the resulting source-native noise parameters are translated into a common pair $(\sigma_A,\tau_c)$. For the Bauch material benchmark, the conversion additionally uses the standard OU echo exponent $\chi_E(t)=\sigma_A^2\tau_c^2[t/\tau_c-3+4e^{-t/(2\tau_c)}-e^{-t/\tau_c}]$~\cite{cywinski2008dephasing}; this equation is used only for parameter inference and does not insert an echo pulse into the simulated gate.

We use the OU model in two complementary ways. First, a Cartesian sensitivity map resolves the dependence on the independently varied Ramsey and correlation times,
\begin{equation}
T_{2,A}^{\ast}\in\{2,5,10,20,50\}\,\mu\mathrm{s},
\qquad
\tau_c\in\{0.1,0.3,1,3,15,30\}\,\mu\mathrm{s}.
\label{eq:ou_cartesian_grid}
\end{equation}
This grid is not interpreted as thirty independently observed material samples. It spans fast noise, gate-timescale correlations, and the approach to the quasi-static regime, while the Ramsey axis resolves the weak-noise limit in which the stochastic curves approach the noiseless control floor.

Second, we evaluate three source-consistent room-temperature material settings as explicit points on the same fidelity plots. For the purified-$^{12}\mathrm C$ CVD material characterized by Bar-Gill \emph{et al.}, we use the central Lorentzian bath parameters $\Delta=30\,\mathrm{kHz}$ and $\tau_c=10\,\mu\mathrm{s}$~\cite{bargill2012spectroscopy}, corresponding in our OU convention to $T_{2,A}^{\ast}=8.57\,\mu\mathrm{s}$. For sample No.~8 of Hayashi \emph{et al.}, the rounded same-sample fit $\lambda\simeq0.025\,\mathrm{MHz}$ and $\tau_c\simeq14\,\mu\mathrm{s}$ gives $\sigma_A/2\pi\simeq50\,\mathrm{kHz}$ and $T_{2,A}^{\ast}=4.76\,\mu\mathrm{s}$~\cite{hayashi2020noise}. Finally, the concentration laws measured by Bauch \emph{et al.} are evaluated at $[\mathrm N]=2\,\mathrm{ppm}$, giving $T_{2,A}^{\ast}=4.8\,\mu\mathrm{s}$ and an OU-equivalent $\tau_c=2.64\,\mathrm{ms}$~\cite{bauch2020decoherence}. The Bauch point is an engineered-material benchmark derived from measured concentration scaling rather than a separately tabulated sample. Source conventions and conversion details are given in Appendix~\ref{app:material_parameter_conversion}.

For every grid point and experimental scenario, the reference and exposure-penalized pulses are propagated using common random numbers. This paired design separates the generic flattening toward the noiseless baseline at large $T_{2,A}^{\ast}$ from robustness created by the control optimization itself: the latter is established only when the exposure-penalized pulse exhibits a smaller stochastic fidelity loss than the reference pulse under the same noise process.

For realization $r$, let $V$ embed the eight-dimensional logical subspace into the full Hilbert space, and let $U_{\mathrm{loc}}$ be the fixed local correction calibrated from the noiseless pulse. The realization-dependent projected logical operator is
\begin{equation}
K_r
=
U_{\mathrm{loc}}V^\dagger U_r(T)V.
\label{eq:noisy_logical_kraus}
\end{equation}
The ensemble map is $\mathcal E(\rho)=N_R^{-1}\sum_rK_r\rho K_r^\dagger$. Projection discards leaked population, so this map is completely positive but generally trace-nonincreasing. For target unitary $U_\star$ and logical dimension $d=8$, we evaluate
\begin{equation}
F_{\mathrm e}
=
\frac{1}{N_Rd^2}
\sum_{r=1}^{N_R}
\left|\operatorname{Tr}(U_\star^\dagger K_r)\right|^2,
\label{eq:ensemble_entanglement_fidelity}
\end{equation}
and the mean logical survival
\begin{equation}
\bar p_{\mathrm{surv}}
=
\frac{1}{N_Rd}
\sum_{r=1}^{N_R}
\operatorname{Tr}(K_r^\dagger K_r).
\label{eq:ensemble_survival}
\end{equation}
Using the Haar-average identity underlying the usual relation between entanglement and average gate fidelity~\cite{nielsen2002average}, while retaining the trace loss associated with leakage, gives the unconditional logical average fidelity
\begin{equation}
F_{\mathrm{avg}}^{(\mathrm{uncond})}
=
\frac{dF_{\mathrm e}+\bar p_{\mathrm{surv}}}{d+1}.
\label{eq:ensemble_average_gate_fidelity}
\end{equation}
Keeping $U_{\mathrm{loc}}$ fixed is essential: fitting a separate correction for every realization would assume knowledge of the instantaneous noise trajectory and overestimate experimentally accessible performance.

\section{Results}
\label{sec:results}

\subsection{Reference tripartite gates and robustness-aware refinement}
\label{subsec:tripartite_gate_synthesis}

The reference controls are the previously obtained direct Fourier pulses for the two targets in Eq.~\eqref{eq:introduction_target_gates}. Both use $N_c=11$ components. The $ZZZ$ reference pulse has duration $1.50\,\mu\mathrm{s}$ and corrected fidelity $0.9984$, while the $XZZ$ reference pulse has duration $1.25\,\mu\mathrm{s}$ and corrected fidelity $0.9985$. The $ZZZ$ target is optimized in the computational frame and the $XZZ$ target in the local Hadamard frame $R=H_A\otimes I_B\otimes I_C$.

Starting from these checkpoints, we perform a continuation in the electron-exposure weight $\lambda_{\mathrm e}$. Each stage is warm-started from the preceding pulse, and the final pulse is selected by minimizing $\mathcal J_{\mathrm e}$ subject to a corrected-fidelity threshold. The pulse duration, $11$-component Fourier family, amplitude bound, support-selective target, and fixed-local-correction convention are unchanged. This isolates the benefit obtained from the open local coefficients rather than from a larger control parametrization or a longer gate.

\bb{Insert the selected continuation weights, final noiseless fidelities, and fractional reductions in $\mathcal J_{\mathrm e}$ and $\overline I_p$ after the production optimization.}

\begin{figure*}[t]
\centering
\includegraphics[width=0.48\textwidth]{Bilder/electron_population_integral_ZZZ.png}
\hfill
\includegraphics[width=0.48\textwidth]{Bilder/electron_population_integral_XZZ.png}
\caption{Electron-trajectory diagnostics for the reference and exposure-penalized controls. Left: $ZZZ$. Right: $XZZ$. The upper panels show the logical-basis-averaged population in the electron manifold opposite to that of each input basis state, with the range over the eight inputs shaded. The lower panels show the mean electronic dephasing exposure $1-\langle\widetilde Z_A\rangle^2$ used in Eq.~\eqref{eq:electron_dephasing_exposure}. The annotated integrals quantify the trajectory change; robustness is independently tested by the OU ensembles below.}
\label{fig:electron_population_integrals}
\end{figure*}

The first-order terminal coefficients are removed using Eq.~\eqref{eq:fixed_local_correction_frame}. For $U_{ZZZ}$ the correction consists of local $Z$ rotations. For $U_{XZZ}$, transforming back from the Hadamard frame gives
\begin{equation}
U_{\mathrm{loc},XZZ}
=
\exp\!\left[-i\left(\phi(A)X_A+\phi(B)Z_B+\phi(C)Z_C\right)\right].
\label{eq:xzz_computational_local_correction}
\end{equation}
The correction is fitted only to the noiseless pulse and is not adjusted for individual noise realizations.

\subsection{Cartesian OU sensitivity maps and experimental scenarios}
\label{subsec:noise_robustness}

For each gate we evaluate both the reference and exposure-penalized pulses on the $5\times6$ Cartesian grid in Eq.~\eqref{eq:ou_cartesian_grid}. The full comparison therefore contains $2$ gates, $2$ pulse variants, and $30$ noise conditions. Common OU trajectories are used for paired pulse comparisons at every point. The plots report unconditional logical average infidelity, including leakage, together with ensemble standard errors and realization quantiles.

\begin{figure*}[t]
\centering
\includegraphics[width=\textwidth]{Bilder/ou_fidelity_sweep_ZZZ_reference_vs_robust.png}
\caption{$ZZZ$ Ramsey-calibrated OU sensitivity map. The left panel shows the reference pulse and the right panel the exposure-penalized pulse. Curves correspond to fixed correlation times and approach the independently evaluated noiseless baseline as $T_{2,A}^{\ast}$ increases. The Bar-Gill/Walsworth, Hayashi, and Bauch/Walsworth material settings are evaluated separately and overlaid as source-consistent markers rather than snapped to the Cartesian grid.}
\label{fig:ou_sweep_zzz}
\end{figure*}

\begin{figure*}[t]
\centering
\includegraphics[width=\textwidth]{Bilder/ou_fidelity_sweep_XZZ_reference_vs_robust.png}
\caption{$XZZ$ Ramsey-calibrated OU sensitivity map, with the same conventions as Fig.~\ref{fig:ou_sweep_zzz}. The paired panels distinguish the generic weak-noise flattening from a reduction in stochastic loss caused by the exposure-aware optimization.}
\label{fig:ou_sweep_xzz}
\end{figure*}

At increasing Ramsey time, every fixed-$\tau_c$ curve must ultimately flatten toward the noiseless control floor because the OU width calibrated by Eq.~\eqref{eq:ou_sigma_ramsey} vanishes. This limiting behavior by itself is not evidence of a robust pulse. The control claim is instead based on the paired displacement between the reference and exposure-penalized panels at identical $(T_{2,A}^{\ast},\tau_c)$, together with the reduced trajectory integrals in Fig.~\ref{fig:electron_population_integrals}. The experimental markers then determine whether this improvement is already relevant in independently characterized room-temperature diamond environments.

\IfFileExists{generated/experimental_scenario_results.tex}{%
\begin{table*}[t]
\centering
\caption{Average gate fidelities in the three experimentally anchored material settings. Stage-one and population-refined pulses use identical noise seeds; uncertainties are one standard error.}
\label{tab:experimental_material_results}
\small
\begin{tabular}{lccccc}
\toprule
Material & $(T_2^*,\tau_c)$ ($\mu$s) & $F_{ZZZ}^{\rm local}$ & $F_{ZZZ}^{\rm pop}$ & $F_{XZZ}^{\rm local}$ & $F_{XZZ}^{\rm pop}$ \\
\midrule
Bar-Gill purified $^{12}\mathrm C$ & $(8.57,10)$ & 0.992320 $\pm$ 1.2e-03 & 0.993211 $\pm$ 1.0e-03 & 0.992904 $\pm$ 1.8e-03 & 0.992921 $\pm$ 1.8e-03 \\
Hayashi HPHT No.~8 & $(4.76,14)$ & 0.985409 $\pm$ 3.2e-03 & 0.986519 $\pm$ 2.9e-03 & 0.980908 $\pm$ 5.1e-03 & 0.980960 $\pm$ 5.1e-03 \\
Bauch engineered $^{12}\mathrm C$ & $(4.8,2.64e+03)$ & 0.984551 $\pm$ 4.0e-03 & 0.986052 $\pm$ 3.2e-03 & 0.980783 $\pm$ 5.7e-03 & 0.980834 $\pm$ 5.7e-03 \\
\bottomrule
\end{tabular}
\end{table*}

}{%
\begin{table*}[t]
\centering
\caption{Placeholder for the automatically generated experimental-scenario results. Run the smoke or production aggregation script to create \texttt{generated/experimental\_scenario\_results.tex}.}
\label{tab:experimental_material_results}
\begin{tabular}{lc}
\toprule
Output & Status \\
\midrule
Bar-Gill, Hayashi, and Bauch scenarios for both gates & Production values pending \\
\bottomrule
\end{tabular}
\end{table*}
}

\bb{After the production run, state quantitatively whether the fidelity loss is minimal or moderate in each experimental scenario and report the paired reduction relative to the reference pulse. Avoid attributing robustness to the large-$T_2^*$ flattening alone.}

\section{Conclusion and Outlook}
\label{sec:conclusion}

We developed a robustness-aware extension of support-selective single-pulse synthesis for tripartite interactions in an NV-center spin register. The phase-coordinate construction can be viewed as a direct many-body generalization of local-equivalence and perfect-entangler optimization: instead of fixing one complete matrix representative, it fixes selected Pauli-support coefficients and leaves correctable local coefficients free. The remaining freedom is then used to reduce the time-integrated electronic $Z$-variance while preserving the required three-body phase and suppressing pairwise content.

The numerical workflow is deliberately paired. The original high-fidelity $ZZZ$ and $XZZ$ pulses serve as warm starts and reference controls. Exposure-penalized pulses are selected at fixed duration and within the same $11$-component Fourier family. Their electron-manifold excursions are diagnosed through time-resolved population and variance integrals, while their actual robustness is determined through stochastic propagation with the OU detuning term included in every short-time propagator.

The Cartesian map over $T_{2,A}^{\ast}\in\{2,5,10,20,50\}\,\mu\mathrm{s}$ and $\tau_c\in\{0.1,0.3,1,3,15,30\}\,\mu\mathrm{s}$ separates fast, gate-timescale, and slowly varying noise. The flattening of each trajectory toward the noiseless floor at large Ramsey time is the expected weak-noise limit. Evidence for control-induced robustness instead comes from the reduced stochastic loss of the exposure-penalized pulse relative to the reference at the same grid points. The Bar-Gill/Walsworth, Hayashi, and Bauch/Walsworth environments place this sensitivity map in experimentally characterized material regimes.

\bb{Replace this paragraph with the final numerical conclusion: selected robust fidelities, exposure reductions, grid-wide improvement, and fidelity loss in the three experimental scenarios.} Subject to that production comparison, the approach provides a concrete route for exchanging otherwise unconstrained local dynamics for reduced exposure of the noisy electronic actuator. Future work includes direct filter-function objectives, robustness to uncertain hyperfine parameters and nuclear noise, and closed-loop refinement on an experimental register.

\section*{Acknowledgements}
\noindent
The authors acknowledge fruitful discussions with Thomas Reisser, Marcus Doherty and Robert Zeier. This work was supported by the European Union’s HORIZON Europe program via project SPINUS (No. 101135699), project QCFD (101080085, HORIZON-CL4-2021 DIGITAL-EMERGING02-10), project OpenSuperQ-Plus100 (101113946, HORIZON-CL4-2022-QUANTUM-01-SGA), by AIDAS-AI, Data Analytics and Scalable Simulation, which is a Joint Virtual Laboratory gathering the Forschungszentrum Juelich and the French Alternative Energies and Atomic Energy Commission, and by Germany’s Excellence Strategy - the Cluster of Excellence Matter and Light for Quantum Computing (ML4Q2) EXC 2004/2–390534769.

\noindent
The authors declare no competing financial interest.

\section*{Author Contributions}
\noindent
B.B. developed the support-selective phase cost used for the NV-center gates, implemented the numerical simulations, optimized the microwave pulses, performed the phase-coordinate analysis, and developed the Ornstein--Uhlenbeck ensemble simulations. Discussions with M.M.M. and F.M. informed the effective NV-register model and its application to tripartite-gate synthesis. M.M.M. and F.M. supervised the project and provided conceptual guidance on the control strategy, noise modeling, and interpretation of the results. T.C., M.M.M., and F.M. contributed scientific guidance, critical feedback, and discussion of the broader implications for multi-spin quantum control. B.B. drafted the manuscript, and all authors reviewed and approved the final version.

\section*{Data Availability}
\noindent
The data and reproduction code are openly available on Zenodo and GitHub~\cite{baran2026_multiqubit_control}.

\appendix

\section{Full Derivation of the NV-Center Hamiltonian}
\label{app:nv_model}

This appendix details the derivation of the effective NV-center Hamiltonian given in Sec.\ref{sec:spin_model}, inspired by the perturbative analysis of Chen \emph{et al.}~\cite{Chen}.


\paragraph{Laboratory Frame.}
Throughout, Hamiltonian parameters are expressed in angular-frequency units. With a static magnetic field $B_0$ aligned to the NV axis and a linearly polarized microwave (MW) field $B_{1,\mathrm{MW}}(t)$ applied orthogonally, the coupled NV-nuclear system is described by
\begin{equation}
\begin{aligned}
H_{\mathrm{NV}}
&= D\!\left(S_z^2-\tfrac{2}{3}\right)
+ \gamma_e B_0 S_z
+ \gamma_e B_{1,\mathrm{MW}}(t)\,S_x\\
&\hspace{-0.7cm}+ \sum_i \vec{S}\!\cdot\! A_i \!\cdot\! \vec{I}_i 
\quad - \sum_i \gamma_i B_0 I_{iz}
\;+\; \sum_i Q_i\!\left(I_{iz}^2-\tfrac{2}{3}\right),
\end{aligned}
\label{eq:H_full}
\end{equation}
where $D/2\pi=2.87\,\mathrm{GHz}$ is the zero-field splitting, $\vec S$ ($S=1$) are the NV electronic spin operators, $\vec I_i$ are nuclear spins ($^{14}$N, selected $^{13}$C) with gyromagnetic ratios $\gamma_i$, $A_i$ are hyperfine tensors, and $Q_i$ are nuclear quadrupole terms.\\

\paragraph{Two-level reduction and secular approximation.}
We consider the effect of applying a strong magnetic field $B_0$. Together with the zero-field splitting $D$, a strong magnetic field lifts the degeneracy of the $m_s\pm1$ electronic levels and isolates the $\{\ket{0_e},\ket{-1_e}\}$ transition from the remaining electronic level $\ket{+1_e}$, enough for it to be neglected (valid for fields of order $0.5\,$T in the parameter regime considered here). The electronic spin can thereby be treated as an effective two-level transition between $\{\ket{0_e},\ket{-1_e}\}$, with transition frequency
\[
\Lambda_s = D - \gamma_e B_0.
\]
With that, we write the electronic Hamiltonian as
\begin{equation}
H_e = \frac{\Lambda_s}{2}\,\sigma_z \;+\; \gamma_e B_{1,\mathrm{MW}}(t)\,\frac{\sigma_x}{\sqrt{2}}.
\label{eq:He_def}
\end{equation}

Given that the electronic splitting $\Lambda_s$ is much larger than the hyperfine couplings $A_i$, the terms in the hyperfine interaction \(\sum_i \vec S\cdot A_i\cdot \vec I_i\) that are non-diagonal in the electronic basis, average out on relevant timescales. This motivates a so-called secular-approximation: only the electron-diagonal components of the hyperfine interaction are retained, so that
\[
\sum_i \vec S\cdot A_i\cdot \vec I_i \longrightarrow
S_z\left(A_{i,zz}I_{iz}+
A_{i,zx}I_{ix}+
A_{i,zy}I_{iy}\right)
\]
Further, we have to consider that in the $\ket{-1_e}$-manifold, the projected hyperfine interaction contributes to the nuclear dynamics with negative sign, giving
\begin{equation}
\begin{aligned}
H_{-1} &= -\sum_i \gamma_i B_0 I_{iz} \;+\; \sum_i Q_i\!\left(I_{iz}^2-\tfrac{2}{3}\right)\\
&\hspace{0.43cm}\;-\; \sum_i \big(A_{i,zz} I_{iz}+A_{i,zx} I_{ix}+A_{i,zy} I_{iy}\big) \label{eq:H-1_def}
\end{aligned}
\end{equation}
By contrast, in the $\ket{0_e}$-manifold the projected $S_z$ contribution vanishes, since $S_z\ket{0_e}=0$, so that the nuclei evolve only under their Zeeman $\gamma_i B_0 I_{iz}$ and quadrupolar terms $Q_i\!\left(I_{iz}^2-\tfrac{2}{3}\right)$,
\begin{equation}
\begin{aligned}
H_{0} &= -\sum_i \gamma_i B_0 I_{iz} \;+\; \sum_i Q_i\!\left(I_{iz}^2-\tfrac{2}{3}\right)
\label{eq:H0_def}
\end{aligned}
\end{equation}
Recombining the secular electronic Hamiltonian with the two conditional nuclear Hamiltonians yields the block form
\begin{equation}
H_{\mathrm{NV}}
= H_e
+ \ket{0_e}\!\bra{0_e}\otimes H_{0}
+ \ket{-1_e}\!\bra{-1_e}\otimes H_{-1}.
\label{eq:block_form}
\end{equation}

\paragraph{Diagonalizing the Nuclear Basis.}
We want to diagonalize the $H_{-1}$ Hamiltonian because it contains off-diagonal terms with respect to the electronic quantization axis, in contrast to the $H_{0}$ Hamiltonian.
In the $H_{-1}$ Hamiltonian, the transverse hyperfine components tilt the effective nuclear field away from the NV axis. For nucleus $i$, the combined secular hyperfine and nuclear Zeeman contributions in the $H_{-1}$ Hamiltonian can be grouped as
\begin{equation}
    H^{(i)}_{-1,\mathrm{hf}} = (A_{i,zz}+\gamma_iB_0)I_{iz}
    +A_{i,zx}I_{ix}
    +A_{i,zy}I_{iy}
    \label{eq:Hminus_hf}
\end{equation}
which, by itself, can be geometrically interpreted as an interaction of the nuclear spin with an effective static field
\[
\vec b_i = 
\begin{pmatrix}
A_{i,zx}\\
A_{i,zy}\\
A_{i,zz}+\gamma_iB_0\\
\end{pmatrix}
\]
whose orientation corresponds to the current quantization axis in the $H_{-1}$ Hamiltonian.


In order to ultimately rotate the basis of $H_{-1}$, we parameterize the orientation of $\vec b_i$ by spherical angles $\theta_i,\phi_i$ \cite{schweiger2001principles},
\begin{equation}
\tan\phi_i = \frac{A_{i,zy}}{A_{i,zx}},
\qquad
\tan\theta_i = \frac{\sqrt{A_{i,zx}^2+A_{i,zy}^2}}
{A_{i,zz}+\gamma_iB_0},
\label{eq:angles}
\end{equation}
and define its magnitude
\begin{equation}
\omega_i = |\vec b_i| =
\sqrt{ (A_{i,zz}+\gamma_iB_0)^2
+ A_{i,zx}^2
+ A_{i,zy}^2 }.
\label{eq:omega_i}
\end{equation}

Based on these spherical angles, we introduce a rotated nuclear basis $(x'_i,y'_i,z'_i)$ whose quantization axis $z'_i$ is aligned with $\vec b_i$. The $I_{iz}$ operator in this rotated basis becomes 
\begin{equation}
    I'_{iz} = I_{ix}\sin\theta_i\cos\phi_i 
    + I_{iy}\sin\theta_i\sin\phi_i 
    + I_{iz}\cos\theta_i.
    \label{eq:Izprime}
\end{equation}
With this choice, the transverse contribution in the $H_{-1}$ Hamiltonian becomes diagonal:
\[
H^{(i)}_{-1,\mathrm{hf}}= \vec b_i\cdot\vec I_i = \omega_i I'_{iz}.
\]

The $H_{-1}$ Hamiltonian therefore takes the form
\begin{align}
H_{-1}
&=
-\sum_i\omega_i I'_{iz}
+
\sum_iQ_i\!\left(I_{iz}'^{\,2}-\tfrac23\right).
\label{eq:H-1_rot}
\end{align}

Thus, after rotating the nuclear basis, the quantization axes associated with the $H_{-1}$ Hamiltonian are collected onto the effective hyperfine axis of each nucleus, rendering the Hamiltonian diagonal in the rotated basis.\\

\paragraph{Electronic transitions between nuclear quantization frames.}

While the $H_0$ Hamiltonian remains quantized along the NV axis, the $H_{-1}$ Hamiltonian is now quantized along the effective hyperfine axis defined by $\vec b_i$. The two nuclear Hamiltonians therefore have a relative tilt between their quantization frames.

The rotated operator \(I'_{iz}\) implicitly defines a nuclear rotation
\(R_i\) satisfying
\[
I'_{iz}=R_i I_{iz} R_i^\dagger.
\]
Since the effective field \(\vec b_i\) is parameterized by the spherical angles \((\theta_i,\phi_i)\), this rotation is generated by
\[
R_i
=
\exp\!\left[
-i\theta_i
\left(
-I_{ix}\sin\phi_i
+
I_{iy}\cos\phi_i
\right)
\right],
\]
with collective rotation
\[
R=\prod_iR_i.
\]

For the weakly misaligned nuclei considered here,  \(\theta_i\ll1\), the rotation may be expanded to first order,
\[
R
\simeq
1-iG,
\]
where
\[
G
=
\sum_i\theta_i
\left(
-I_{ix}\sin\phi_i
+
I_{iy}\cos\phi_i
\right).
\]

Rather than rewriting both nuclear Hamiltonians in a common operator basis, we absorb this relative nuclear rotation into the electronic transition operators. Defining the electronic raising and lowering operators as
\[
\sigma_+
=
\ket{0_e}\!\bra{-1_e},
\qquad
\sigma_-
=
\ket{-1_e}\!\bra{0_e},
\]
the microwave interaction
\[
H_{\mathrm{MW}}
=
\frac{\gamma_eB_{1,\mathrm{MW}}(t)}{\sqrt2}
(\sigma_+ + \sigma_-)
\]
therefore becomes
\begin{equation}
\begin{aligned}
H_{\mathrm{MW}}
&=
\frac{\gamma_eB_{1,\mathrm{MW}}(t)}{\sqrt2}
\left(
\sigma_+\otimes R^\dagger
+
\sigma_-\otimes R
\right)
\\
&\approx
\frac{\gamma_eB_{1,\mathrm{MW}}(t)}{\sqrt2}
\left[
\sigma_x
-
\sigma_y G
\right].
\end{aligned}
\label{eq:HMW_rot_first_order}
\end{equation}

The first term describes the bare electronic microwave transition, while the second term is a first-order electron-nuclear modulation arising from the mismatch between the nuclear quantization frames of the two electronic manifolds.\\


\paragraph{Interaction picture.}

The next step of the derivation is a transformation to the interaction picture that removes the free precession dynamics and leaves only the control-induced dynamics explicit. We therefore decompose the Hamiltonian as
\begin{equation}
H_{\mathrm{NV}}
=
H_{\mathrm{free}}
+
H_{\mathrm{MW}},
\label{eq:H_split}
\end{equation}
where the free evolution is generated by the diagonal electronic and nuclear
terms,
\begin{equation}
\begin{aligned}
H_{\mathrm{free}}
&=
\frac{\Lambda_s}{2}\sigma_z
\\
&\quad
+
\ket{0_e}\!\bra{0_e}\otimes
\left[
-\sum_i\gamma_iB_0 I_{iz}
+
\sum_iQ_i\!\left(I_{iz}^{\,2}-\tfrac23\right)
\right]
\\
&\quad
+
\ket{-1_e}\!\bra{-1_e}\otimes
\left[
-\sum_i\omega_i I'_{iz}
+
\sum_iQ_i\!\left(I_{iz}'^{\,2}-\tfrac23\right)
\right],
\end{aligned}
\label{eq:H_free}
\end{equation}
while the microwave interaction is
\begin{equation}
\begin{aligned}
H_{\mathrm{MW}}
&\approx
\frac{\gamma_e B_{1,\mathrm{MW}}(t)}{\sqrt2}
\left[
\sigma_x
-
\sigma_y G
\right]\,.
\end{aligned}
\label{eq:HMW_rot}
\end{equation}


We choose \(H_{\mathrm{free}}\) as the generator of the interaction picture, such that
\begin{equation}
H_{\mathrm{NV}}^{(I)}(t)
=
e^{iH_{\mathrm{free}}t}
H_{\mathrm{MW}}
e^{-iH_{\mathrm{free}}t}.
\label{eq:IP_transform}
\end{equation}

We first transform the bare electronic part of the microwave interaction,
\[
H_{\mathrm{MW},x}
=
\frac{\gamma_e B_{1,\mathrm{MW}}(t)}{\sqrt2}\sigma_x .
\]
Its interaction-picture form is
\begin{equation}
H_{\mathrm{MW},x}^{(I)}(t)
=
\frac{\gamma_e B_{1,\mathrm{MW}}(t)}{\sqrt2}
e^{iH_{\mathrm{free}}t}
\sigma_x
e^{-iH_{\mathrm{free}}t}.
\label{eq:HMWx_IP_start}
\end{equation}


Since \(H_{\mathrm{free}}\) is diagonal in the electronic basis and in the natural nuclear basis of each manifold, its propagator acts only through energy-dependent phase accumulation,
\[
e^{iH_{\mathrm{free}}t}
=
\sum_{e,\mathbf m}
e^{iE_e(\mathbf m)t}
\ket{e,\mathbf m}\!\bra{e,\mathbf m}.
\]
Using \(\sigma_x=\sigma_++\sigma_-\), the interaction picture evolution of the electronic transition operator is therefore governed by the energy difference between the connected free eigenstates,
\begin{equation}
\begin{aligned}
&e^{iH_{\mathrm{free}}t}
\sigma_x
e^{-iH_{\mathrm{free}}t}
\notag\\
&=
\sum_{\mathbf m}
\Big[
\sigma_+
e^{i[E_0(\mathbf m)-E_{-1}(\mathbf m)+\Lambda_s]t}
\\
&\hspace{1.5cm}
+
\sigma_-
e^{-i[E_0(\mathbf m)-E_{-1}(\mathbf m)+\Lambda_s]t}
\Big]
\otimes
\ket{\mathbf m}\!\bra{\mathbf m}.
\end{aligned}
\label{eq:sigmax_IP_phase}
\end{equation}
Defining the configuration-dependent transition frequency
\begin{equation}
\begin{aligned}
\Lambda(\mathbf m)
&:=
\Lambda_s
+
E_0(\mathbf m)
-
E_{-1}(\mathbf m)
\\
&=
\Lambda_s
+
\left[
-\sum_i\gamma_iB_0m_i
+
\sum_iQ_i\!\left(m_i^2-\tfrac23\right)
\right]
\\
&\quad
-
\left[
-\sum_i\omega_i m_i
+
\sum_iQ_i\!\left(m_i^2-\tfrac23\right)
\right]
\\
&=
\Lambda_s
-
\sum_i m_i(\gamma_iB_0-\omega_i).
\end{aligned}
\label{eq:Lambda_config_general}
\end{equation}
this becomes
\begin{equation}
\begin{aligned}
&e^{iH_{\mathrm{free}}t}
\sigma_x
e^{-iH_{\mathrm{free}}t}
\notag\\
&=
\sum_{\mathbf m}
\left[
\sigma_+ e^{i\Lambda(\mathbf m)t}
+
\sigma_- e^{-i\Lambda(\mathbf m)t}
\right]
\otimes
\ket{\mathbf m}\!\bra{\mathbf m}
\\
&=
\sum_{\mathbf m}
\big[
\sigma_x\cos(\Lambda(\mathbf m)t)
-
\sigma_y\sin(\Lambda(\mathbf m)t)
\big]
\otimes
\ket{\mathbf m}\!\bra{\mathbf m}.
\end{aligned}
\label{eq:sigmax_IP_step}
\end{equation}


Consequently, the bare electronic part of the microwave interaction becomes
\begin{equation}
\begin{aligned}
H_{\mathrm{MW}}^{(I)}(t)
&=
\frac{\gamma_e B_{1,\mathrm{MW}}(t)}{\sqrt2}
\sum_{\mathbf m}
\big[
\sigma_x\cos(\Lambda(\mathbf m)t)
-
\sigma_y\sin(\Lambda(\mathbf m)t)
\big]
\\
&\hspace{1cm}\otimes
\ket{\mathbf m}\!\bra{\mathbf m}.
\end{aligned}
\label{eq:HMW_I}
\end{equation}
Here, the nuclear configuration is preserved to leading order during the bare
microwave-driven electronic transition.\\

%


We now transform the second term in the dressed microwave interaction,
\[
H_{\mathrm{MW},eN}
=
-\frac{\gamma_e B_{1,\mathrm{MW}}(t)}{\sqrt2}
\sigma_y\otimes G .
\]
Its interaction-picture form is
\begin{equation}
H_{\mathrm{MW},eN}^{(I)}(t)
=
-\frac{\gamma_e B_{1,\mathrm{MW}}(t)}{\sqrt2}
e^{iH_{\mathrm{free}}t}
\left(
\sigma_y\otimes G
\right)
e^{-iH_{\mathrm{free}}t}.
\label{eq:HeN_IP_start}
\end{equation}

Since the electron-nuclear modulation term contributes only at first order in the small misalignment angles \(\theta_i\), we employ a leading-order description of its interaction-picture evolution. For the electronic part, we retain the dominant bare precession at \(\Lambda_s\) rather than the full configuration-dependent transition frequency \(\Lambda(\mathbf m)\). For the nuclear part, we approximate the manifold-dependent precession by the effective mean generator
\[
H_N
=
-\sum_i \delta_i I_{iz},
\,\, \text{with}\,\,
\delta_i
=
\frac{\gamma_iB_0+\omega_i}{2}.
\]
This captures the leading averaged dynamics of the weak modulation without
tracking the full configuration-resolved evolution.

For the same reason, we also approximate the transformation as
\begin{equation}
\begin{aligned}
&e^{iH_{\mathrm{free}}t}
\left(
\sigma_y\otimes G
\right)
e^{-iH_{\mathrm{free}}t}
\notag\\
&\approx
\left(
e^{i\Lambda_s\sigma_z t/2}
\sigma_y
e^{-i\Lambda_s\sigma_z t/2}
\right)
\otimes
\left(
e^{iH_N t}
G
e^{-iH_N t}
\right).
\end{aligned}
\label{eq:sigmayG_factorized}
\end{equation}

The electronic part evaluates to
\[
e^{i\Lambda_s\sigma_z t/2}
\sigma_y
e^{-i\Lambda_s\sigma_z t/2}
=
\sigma_x\sin(\Lambda_s t)
+
\sigma_y\cos(\Lambda_s t),
\]
while the nuclear operator
\[
G
=
\sum_i\theta_i
\left(
-I_{ix}\sin\phi_i
+
I_{iy}\cos\phi_i
\right)
\]
evolves under \(H_N\) as
\[
e^{iH_N t}
G
e^{-iH_N t}
=
\sum_i\theta_i
\left[
I_{ix}\cos(\delta_i t-\phi_i)
+
I_{iy}\sin(\delta_i t-\phi_i)
\right].
\]

Therefore, to leading order,
\begin{equation}
\begin{aligned}
H_{\mathrm{MW},eN}^{(I)}(t)
&\approx
-\frac{\gamma_e B_{1,\mathrm{MW}}(t)}{\sqrt2}
\big[
\sigma_x\sin(\Lambda_s t)
+
\sigma_y\cos(\Lambda_s t)
\big]
\\
&\quad\otimes
\sum_i\theta_i
\big[
I_{ix}\cos(\delta_i t-\phi_i)
+
I_{iy}\sin(\delta_i t-\phi_i)
\big].
\end{aligned}
\label{eq:HeN_I}
\end{equation}


\paragraph{Rotating-wave approximation (RWA).}
We drive the electron at microwave frequency $\omega_{\mathrm{MW}}$ with a magnetic field
\begin{equation}
B_{1,\mathrm{MW}}(t)
=
\frac{\sqrt2\,\Omega(t)}{\gamma_e}
\cos(\omega_{\mathrm{MW}} t),
\end{equation}
where $\Omega(t)$ denotes the effective time-dependent Rabi-frequency envelope. 
Substituting this field into Eq.~\eqref{eq:HMW_I} gives the prefactor
\[
\frac{\gamma_e B_{1,\mathrm{MW}}(t)}{\sqrt2}
=
\Omega(t)\cos(\omega_{\mathrm{MW}} t).
\]
Using the product-to-sum identity $\cos\alpha\cos\beta=\tfrac12[\cos(\alpha-\beta)+\cos(\alpha+\beta)]$
and $\cos\alpha\sin\beta=\tfrac12[\sin(\beta+\alpha)+\sin(\beta-\alpha)]$,
the electron interaction-picture drive
\(
\propto \sigma_x\cos[\Lambda(\mathbf m)t]-\sigma_y\sin[\Lambda(\mathbf m)t]
\)
splits into a slow (difference-frequency) piece and explicit counter-rotating (sum-frequency) pieces:
\begin{align}
&\notag\hspace{-0.0cm}H_{e}^{(\mathrm{full})}(t)= \frac{\Omega(t)}{2}\sum_{\mathbf m}\\
&\hspace{-0.2cm}
\Big\{\underbrace{\big[\sigma_x\cos((\omega_{\mathrm{MW}}-\Lambda(\mathbf m))\,t)-\sigma_y\sin((\omega_{\mathrm{MW}}-\Lambda(\mathbf m))\,t)\big]}_{\text{slow (near-resonant)}}\nonumber\\
&\hspace{-0.2cm}
+\notag \underbrace{\big[\sigma_x\cos\big((\omega_{\mathrm{MW}}+\Lambda(\mathbf m))t\big)- 
\hspace{0.2cm} \sigma_y\sin\big((\omega_{\mathrm{MW}}+\Lambda(\mathbf m))t\big)\big]}_{\text{fast (counter-rotating)}}\Big\}\\
&\otimes \ket{\mathbf m}\!\bra{\mathbf m}.
\label{eq:He_full_split}
\end{align}
The \emph{rotating-wave approximation} (RWA) keeps only the slow terms and discards the counter-rotating term at  $(\omega_{\mathrm{MW}}+\Lambda(\mathbf m))t\big)$, yielding
\begin{equation}
\boxed{
\begin{aligned}
&H_{e}^{(\mathrm{RWA})}(t)
= \frac{\Omega(t)}{2} \sum_{\mathbf m}
\big[\sigma_x\cos((\omega_{\mathrm{MW}}-\Lambda(\mathbf m))\,t)\\ & \hspace{1.9cm}-\sigma_y\sin((\omega_{\mathrm{MW}}-\Lambda(\mathbf m))\,t)\big]
\otimes \ket{\mathbf m}\!\bra{\mathbf m}.
\end{aligned}
}
\label{eq:HeRWA_time_nophase_Lambda_m}
\end{equation}
Applying the same procedure to the electron–nuclear modulation term Eq.~\eqref{eq:HeN_I}, gives
\begin{align}
&\notag H_{eN}^{(\mathrm{full})}(t)\\
&\notag=
\underbrace{-\,\frac{\Omega(t)}{2}\,
\big[\sigma_x\sin((\omega_{\mathrm{MW}}-\Lambda_s) t)+\sigma_y\cos((\omega_{\mathrm{MW}}-\Lambda_s) t)\big]}_{\text{slow (near-resonant in the electron sector)}}\;\\
&\notag\;\otimes\;
\sum_i \theta_i\big[I'_{ix}\cos(\delta_i t-\phi_i)+I'_{iy}\sin(\delta_i t-\phi_i)\big]\\
&\notag\hspace{-0.0cm}\quad
\underbrace{-\,\frac{\Omega(t)}{2}\,
\big[\sigma_x\sin\big((\Lambda_s+\omega_{\mathrm{MW}})t\big)
      +\sigma_y\cos\big((\Lambda_s+\omega_{\mathrm{MW}})t\big)\big]}_{\text{fast (counter-rotating)}}\\
&\;\otimes\;
\sum_i \theta_i\big[I'_{ix}\cos(\delta_i t-\phi_i)+I'_{iy}\sin(\delta_i t-\phi_i)\big].
\label{eq:HeN_full_split}
\end{align}
Discarding the fast rotating term again, gives 
\begin{equation}
\boxed{
\begin{aligned}
&H_{eN}^{(\mathrm{RWA})}(t)\\
&\notag=-\,\frac{\Omega(t)}{2}\,
\big[\sigma_x\sin((\omega_{\mathrm{MW}}-\Lambda_s) t)+\sigma_y\cos((\omega_{\mathrm{MW}}-\Lambda_s) t)\big]\\
&\;\otimes\;
\sum_i \theta_i\big[I'_{ix}\cos(\delta_i t-\phi_i)+I'_{iy}\sin(\delta_i t-\phi_i)\big].
\end{aligned}
}
\label{eq:HeN_RWA_time_nophase_final}
\end{equation}

\section{Conversion of Material Characterizations to OU Parameters}
\label{app:material_parameter_conversion}

The three source experiments characterize their environments using
different conventions. We translate each material setting into the
common autocorrelation
\begin{equation}
\left\langle\beta(t)\beta(0)\right\rangle
=
\sigma_A^2e^{-|t|/\tau_c}
\end{equation}
used in the stochastic propagation.

For the purified-$^{12}\mathrm C$ material of Bar-Gill
\emph{et al.}, we identify the reported central Lorentzian scale
$\Delta=30\,\mathrm{kHz}$ with the rms transition-frequency
fluctuation, giving $\sigma_A/2\pi=30\,\mathrm{kHz}$, and retain the
measured central correlation time $\tau_c=10\,\mu\mathrm{s}$
~\cite{bargill2012spectroscopy}.

Hayashi \emph{et al.} write the stochastic interaction as
\begin{equation}
H_I(t)=\lambda f(t)\sigma_z,
\end{equation}
where $f(t)$ has unit stationary variance. Equating this expression
with $H_{\mathrm{noise}}(t)=\beta(t)Z/2$ gives
$\beta(t)=2\lambda f(t)$. After conversion from the reported
frequency units to angular-frequency units,
\begin{equation}
\frac{\sigma_A}{2\pi}=2\lambda.
\end{equation}
For the rounded sample-No.~8 value
$\lambda\simeq0.025\,\mathrm{MHz}$, this gives
$\sigma_A/2\pi\simeq50\,\mathrm{kHz}$, while
$\tau_c\simeq14\,\mu\mathrm{s}$ is retained from the same-sample fit
~\cite{hayashi2020noise}.

For the Bauch benchmark, the measured concentration laws are
evaluated at $[\mathrm N]=2\,\mathrm{ppm}$, yielding
$T_{2,A}^{\ast}=4.8\,\mu\mathrm{s}$ and
$T_2=71.7\,\mu\mathrm{s}$~\cite{bauch2020decoherence}. The OU
parameters are obtained by solving
\begin{equation}
\chi_R(T_{2,A}^{\ast})=1,
\qquad
\chi_E(T_2)=1,
\end{equation}
where $\chi_R$ is defined in Eq.~\eqref{eq:ou_ramsey_envelope} and
\begin{equation}
\chi_E(t)
=
\sigma_A^2\tau_c^2
\left(
\frac{t}{\tau_c}-3
+4e^{-t/(2\tau_c)}
-e^{-t/\tau_c}
\right).
\end{equation}
The resulting effective parameters are
\begin{equation}
\frac{\sigma_A}{2\pi}=46.9\,\mathrm{kHz},
\qquad
\tau_c=2.64\,\mathrm{ms}.
\end{equation}
The echo expression is used only for parameter inference; no echo
pulse is included in the gate propagation.

\section{Explicit Three-Qubit Phase Coefficients}
\label{app:explicit_3q_invariants}

For three qubits, the support-selective phase coefficients follow directly from Eq.~\eqref{eq:support_phase_coordinate}. Writing $x=(a,b,c)$ and denoting the computational-basis phases by $\varphi_{abc}$ gives
\begin{equation}
\phi(S)
=
\frac{1}{8}
\sum_{a,b,c\in\{0,1\}}
(-1)^{\sum_{j\in S}x_j}\,\varphi_{abc}.
\label{eq:phiS_3q_general}
\end{equation}
The sign of each basis-state phase is fixed by the parity of the bits on the selected support $S$. The seven non-global coefficients are listed explicitly below.

\medskip
\noindent\textbf{Tripartite coefficient.}
\begin{equation}
\boxed{
\begin{aligned}
\phi(ABC)
&=\frac{1}{8}\Big(
\varphi_{000}-\varphi_{001}-\varphi_{010}+\varphi_{011}\\
&\qquad\quad
-\varphi_{100}+\varphi_{101}+\varphi_{110}-\varphi_{111}
\Big).
\end{aligned}
}
\end{equation}

\medskip
\noindent\textbf{Pairwise coefficients.}
\begin{equation}
\boxed{
\begin{aligned}
\phi(AB)
&=\frac{1}{8}\Big(
\varphi_{000}+\varphi_{001}-\varphi_{010}-\varphi_{011}\\
&\qquad\quad
-\varphi_{100}-\varphi_{101}+\varphi_{110}+\varphi_{111}
\Big),
\end{aligned}
}
\end{equation}
\begin{equation}
\boxed{
\begin{aligned}
\phi(AC)
&=\frac{1}{8}\Big(
\varphi_{000}-\varphi_{001}+\varphi_{010}-\varphi_{011}\\
&\qquad\quad
-\varphi_{100}+\varphi_{101}-\varphi_{110}+\varphi_{111}
\Big),
\end{aligned}
}
\end{equation}
\begin{equation}
\boxed{
\begin{aligned}
\phi(BC)
&=\frac{1}{8}\Big(
\varphi_{000}-\varphi_{001}-\varphi_{010}+\varphi_{011}\\
&\qquad\quad
+\varphi_{100}-\varphi_{101}-\varphi_{110}+\varphi_{111}
\Big).
\end{aligned}
}
\end{equation}

\medskip
\noindent\textbf{One-body coefficients.}
\begin{equation}
\boxed{
\begin{aligned}
\phi(A)
&=\frac{1}{8}\Big(
\varphi_{000}+\varphi_{001}+\varphi_{010}+\varphi_{011}\\
&\qquad\quad
-\varphi_{100}-\varphi_{101}-\varphi_{110}-\varphi_{111}
\Big),
\end{aligned}
}
\end{equation}
\begin{equation}
\boxed{
\begin{aligned}
\phi(B)
&=\frac{1}{8}\Big(
\varphi_{000}+\varphi_{001}-\varphi_{010}-\varphi_{011}\\
&\qquad\quad
+\varphi_{100}+\varphi_{101}-\varphi_{110}-\varphi_{111}
\Big),
\end{aligned}
}
\end{equation}
\begin{equation}
\boxed{
\begin{aligned}
\phi(C)
&=\frac{1}{8}\Big(
\varphi_{000}-\varphi_{001}+\varphi_{010}-\varphi_{011}\\
&\qquad\quad
+\varphi_{100}-\varphi_{101}+\varphi_{110}-\varphi_{111}
\Big).
\end{aligned}
}
\end{equation}

\bibliographystyle{unsrt}
\bibliography{references,three_material_environment_references}

\end{document}

